\chapter{Proposed Methodology}
\label{Proposed_Methodology}

This chapter details the implemented information-hiding pipeline, from system architecture and per-modality embedding/extraction routines to evaluation metrics and implementation details. The chapter is organized in a top-down fashion: Section~\ref{sec:architecture} presents the system architecture and a block diagram of the application; Section~\ref{sec:image_module} details the row-based LSB image-watermarking algorithm; Sections~\ref{sec:audio_module}, \ref{sec:text_module} and~\ref{sec:video_module} describe the audio, text and video modules respectively; Section~\ref{sec:payload_format} formalizes the common payload format; Section~\ref{sec:complexity} performs a worst-case complexity analysis; Section~\ref{sec:metrics} specifies the evaluation metrics used in Chapter~\ref{Experimental_results}; Section~\ref{sec:implementation} lists the implementation environment and reproducibility provisions; and Section~\ref{sec:methodology_summary} summarizes the methodology.

\section{System Architecture Overview}
\label{sec:architecture}

The proposed application is implemented in Java SE~17 using Java Swing for the graphical user interface. The architecture keeps the GUI separate from four carrier-specific modules, as shown in Figure~\ref{fig:architecture}.

\begin{figure}[H]
    \centering
    \resizebox{\textwidth}{!}{%
    \begin{tikzpicture}[
        node distance=0.95cm and 1.2cm,
        every node/.style={font=\small},
        layer/.style={draw, rounded corners=3pt, text width=12.8cm, minimum height=0.95cm, align=center, fill=gray!10},
        process/.style={draw, rounded corners=2pt, text width=3.45cm, minimum height=0.9cm, align=center, fill=white},
        io/.style={draw, rounded corners=8pt, text width=3.6cm, minimum height=0.85cm, align=center, fill=white},
        note/.style={font=\scriptsize, align=left, text width=2.35cm},
        arrow/.style={->, >=stealth, semithick}
    ]
        \node[layer] (gui) {\textbf{Swing GUI and Control Layer}\\Watermark input, password input, file loading, action buttons, image preview, and in-window log console};
        \node[layer, below=of gui] (shared) {\textbf{Shared Processing Layer}\\Payload creation with timestamp, logger callback, and per-modality dispatch from the main application};

        % Image branch
        \node[io, below left=1.5cm and 3.4cm of shared] (imgin) {BMP cover image\\($256 \times 256$)};
        \node[process, below=of imgin] (imgp1) {MD5 hash of password\\Row and start-column selection};
        \node[process, below=of imgp1] (imgp2) {Row occupation check\\3--3--2 RGB LSB embedding};
        \node[io, below=of imgp2] (imgout) {Watermarked BMP\\and extracted watermark};
        \node[note, left=0.25cm of imgp1.west, anchor=east] {Password chooses\\embedding row and\\starting column};

        % Audio branch
        \node[io, right=1.35cm of imgin] (audin) {WAV audio file};
        \node[process, below=of audin] (audp1) {Watermark + timestamp\\XOR encryption};
        \node[process, below=of audp1] (audp2) {Bitstream generation\\LSB substitution in samples};
        \node[io, below=of audp2] (audout) {Stego WAV\\and recovered message};
        \node[note, right=0.25cm of audp1.east, anchor=west] {End marker:\\\texttt{1111111111111110}};

        % Text branch
        \node[io, below right=1.55cm and 3.4cm of shared] (txtin) {Plain-text /\\Markdown file};
        \node[process, below=of txtin] (txtp1) {Watermark +\\timestamp\\XOR encryption};
        \node[process, below=of txtp1] (txtp2) {Append encrypted block\\inside \texttt{<!--STEGO:...-->}};
        \node[io, below=of txtp2] (txtout) {Stego text file\\and recovered message};

        % Video branch
        \node[io, left=1.35cm of txtin] (vidin) {Input video file};
        \node[process, below=of vidin] (vidp1) {Extract frame 0 with FFmpeg\\Use top-left $256 \times 256$ region};
        \node[process, below=of vidp1] (vidp2) {Apply row-based image watermark\\Lossless FFV1 re-encoding};
        \node[io, below=of vidp2] (vidout) {Stego MKV\\and recovered watermark};
        \node[note, left=0.2cm of vidp1.west, anchor=east, text width=1.85cm] (vidnote) {Video module\\reuses the\\image algorithm};

        % Shared arrows
        \draw[arrow] (gui) -- (shared);
        \draw[arrow] (shared.south west) .. controls +(-1.0,-0.6) and +(0,0.9) .. (imgin.north);
        \draw[arrow] (shared.south west) .. controls +(0.4,-0.6) and +(0,0.9) .. (audin.north);
        \draw[arrow] (shared.south east) .. controls +(-0.4,-0.6) and +(0,0.9) .. (vidin.north);
        \draw[arrow] (shared.south east) .. controls +(1.0,-0.6) and +(0,0.9) .. (txtin.north);

        % Vertical flows
        \draw[arrow] (imgin) -- (imgp1);
        \draw[arrow] (imgp1) -- (imgp2);
        \draw[arrow] (imgp2) -- (imgout);

        \draw[arrow] (audin) -- (audp1);
        \draw[arrow] (audp1) -- (audp2);
        \draw[arrow] (audp2) -- (audout);

        \draw[arrow] (txtin) -- (txtp1);
        \draw[arrow] (txtp1) -- (txtp2);
        \draw[arrow] (txtp2) -- (txtout);

        \draw[arrow] (vidin) -- (vidp1);
        \draw[arrow] (vidp1) -- (vidp2);
        \draw[arrow] (vidp2) -- (vidout);
    \end{tikzpicture}
    }
    \caption{System architecture of the proposed watermarking framework. A common Swing-based control layer collects the user input and dispatches the request to four carrier-specific pipelines for image, audio, text, and video processing.}
    \label{fig:architecture}
\end{figure}

The architecture consists of four primary components:

\begin{enumerate}
    \item \textbf{GUI Layer (Swing)}: collects the watermark text, the password and the user action (\texttt{Load}, \texttt{Embed}, \texttt{Extract}, \texttt{Save}) from a single window and dispatches it to the appropriate algorithmic module.
    \item \textbf{Per-Modality Algorithmic Modules}: four self-contained classes (\texttt{DWM2}, \texttt{AudioSteganography}, \texttt{TextSteganography}, \texttt{VideoSteganography}) that implement the carrier-specific embedding and extraction routines.
    \item \textbf{Cryptographic Helpers}: an MD5-plus-XOR-fold helper for the image (and video) module, and an XOR-encryption helper for the audio and text modules.
    \item \textbf{Common Payload Layer}: a shared self-delimited timestamped payload format (Section~\ref{sec:payload_format}) that all four modules read and write.
\end{enumerate}

This layered strategy lets each module address its carrier-specific challenges while keeping the user experience and the recovered-payload parser uniform. Each module is implemented as a stand-alone class with no compile-time dependency on the GUI; the GUI passes a logger callback (a \verb|Consumer<String>|) to the module so that the module can write its progress messages back into the in-window log without referencing any Swing class.

\subsection{Typical User Flow}

From a user's point of view, the system is intentionally simple. The user opens the application, selects a file, types a short watermark, enters a password, and chooses whether to embed or extract. The same small set of actions is used for image, audio, text, and video. This reduces confusion because the user does not need to switch to a different tool or learn a different sequence of steps for each media type.

During embedding, the GUI first checks that the required inputs are present. It then forwards the request to the correct module based on the selected file. The module builds the payload, performs the carrier-specific embedding, and sends progress messages back to the GUI log. After the operation finishes, the user can preview the result when relevant and save the stego-file to disk.

During extraction, the flow is similar but reversed. The selected module reads the hidden data from the file, reconstructs the payload, and separates the watermark from the timestamp. The recovered values are then shown to the user in the same application window. This symmetry between embedding and extraction is one of the practical strengths of the design.

\subsection{Why the System is Organized this Way}

The architecture is not complicated, but each design choice has a reason behind it. The GUI layer is kept separate so that interface code does not get mixed with watermarking logic. This makes the source code easier to read and also makes testing simpler, because the core methods can be called without starting the full interface.

The shared processing layer exists because some tasks repeat across multiple modules. For example, all modules need to accept a watermark, use a password, generate a payload, and report status messages. Keeping these ideas conceptually grouped helps the project feel like one framework rather than four unrelated mini-programs.

At the same time, each carrier needs its own embedding rule. Pixels, audio samples, text comments, and video frames behave very differently. Trying to force them into one low-level algorithm would make the code harder to follow and would not improve the result. For this reason, the project uses a common high-level structure with separate low-level implementations.

\section{Image-Watermarking Module}
\label{sec:image_module}

The image-watermarking module is the main part of the proposed framework. It accepts a $256 \times 256$ BMP cover image, a watermark string of at most $16$ ASCII characters, and a non-empty password. It then produces a stego-image that stores the watermark in the LSBs of one row of pixels.

\subsection{Password-to-Location Mapping}
\label{sec:password_mapping}

Let $p$ denote the user-supplied password and let $H = \mathrm{MD5}(p)$ denote its 128-bit MD5 digest, which is represented as a sequence of $16$ bytes $H = (h_0, h_1, \ldots, h_{15})$. Two non-overlapping XOR folds of $H$ produce the row index $r$ and the starting column index $c$:
\begin{align}
    r &= h_0 \oplus h_1 \oplus h_2 \oplus \cdots \oplus h_{11}, \label{eq:row}\\
    c &= h_{12} \oplus h_{13} \oplus h_{14} \oplus h_{15}. \label{eq:col}
\end{align}
Because each $h_i \in \{0, \ldots, 255\}$ and XOR is closed on this set, both $r$ and $c$ lie in $\{0, \ldots, 255\}$. In practical terms, this means the password can spread different watermarks across many possible row and column pairs instead of always placing them in one fixed position. Using separate parts of the digest for $r$ and $c$ also helps the row choice and the starting-column choice behave like two separate decisions.

\subsection{Per-Pixel 3--3--2 Bit Split}

Each pixel in the cover image is a 24-bit RGB triple $(R, G, B)$ with $R, G, B \in \{0, \ldots, 255\}$. To embed an 8-bit ASCII character $\chi$ with binary representation $b_7 b_6 b_5 b_4 b_3 b_2 b_1 b_0$ into a single pixel, the algorithm splits $\chi$ into a 3-bit slice for red, a 3-bit slice for green and a 2-bit slice for blue:
\begin{equation}
    \chi = \underbrace{b_7 b_6 b_5}_{\text{Red LSB}} \mid \underbrace{b_4 b_3 b_2}_{\text{Green LSB}} \mid \underbrace{b_1 b_0}_{\text{Blue LSB}} \enspace .
\end{equation}
The new pixel $(R', G', B')$ is then computed by zeroing the three least-significant bits of $R$ and $G$ and the two least-significant bits of $B$, and substituting in the three slices:
\begin{align}
    R' &= (R \;\&\; \mathtt{0xF8}) \mid (b_7 b_6 b_5)_2, \label{eq:redmod}\\
    G' &= (G \;\&\; \mathtt{0xF8}) \mid (b_4 b_3 b_2)_2, \\
    B' &= (B \;\&\; \mathtt{0xFC}) \mid (b_1 b_0)_2.
\end{align}
The maximum per-channel perturbation introduced by~\eqref{eq:redmod} is therefore $|R' - R| \le 7$, $|G' - G| \le 7$ and $|B' - B| \le 3$, which is well below the just-noticeable-difference threshold of the human visual system (typically $\sim 5$ on each channel for natural images). Reversing the encoding requires only the lower-bit masks:
\begin{equation}
    \chi = ((R' \;\&\; \mathtt{0x07}) \ll 5) \mid ((G' \;\&\; \mathtt{0x07}) \ll 2) \mid (B' \;\&\; \mathtt{0x03}).
\end{equation}

\subsection{Row-Occupation Marker}

A potential collision arises when two different passwords map to the same row $r$ (an event that occurs with probability $1/256$ per pair). To prevent silent overwrites, the embedder writes the sentinel character `\$' (ASCII~36) into pixel $(0, r)$ of every occupied row, using the same 3--3--2 encoding. Before any new embedding the routine reads pixel $(0, r)$, decodes it and checks whether the recovered character is `\$'. If it is, the embedding is refused with an explicit error message that asks the user to choose a different password; otherwise the routine proceeds with the embedding and writes the marker.

\subsection{Self-Delimited Payload Layout}

The payload that is actually written into the row consists of the user watermark $w$ surrounded by three delimiter characters and the embedding timestamp $t$:
\begin{equation}
    \pi(w, t) = \mathtt{*} \;\Vert\; w \;\Vert\; \mathtt{\#} \;\Vert\; t \;\Vert\; \mathtt{\#@} \enspace ,
\end{equation}
where $\Vert$ denotes string concatenation, `\texttt{*}' (ASCII~42) marks the start of the payload, `\texttt{\#}' (ASCII~35) separates the watermark from the timestamp and `\texttt{\#@}' marks the end. The first character of $\pi(w, t)$ is written at column $c$ of row $r$, and the subsequent characters are written at columns $(c+1) \bmod 256, (c+2) \bmod 256, \ldots$, wrapping around within the same row. Because $|w| \le 16$ and the timestamp has fixed length $19$, the total payload length is at most $1 + 16 + 1 + 19 + 2 = 39$ characters, which fits comfortably in the $255$ usable pixels of a row (column $0$ is reserved for the `\$' marker).

\subsection{Embedding Algorithm}

Algorithm~\ref{alg:embed} formalizes the complete image embedding routine.

\begin{algorithm}[H]
    \RaggedRight
    \caption{Row-based LSB watermark embedding}
    \label{alg:embed}
    \algorithmicrequire \quad Cover image $I \in \{0,\ldots,255\}^{256 \times 256 \times 3}$, watermark $w$ with $|w| \le 16$, password $p$. \\
    \algorithmicensure \quad Stego-image $I'$ containing $w$ in row $r$.

    \begin{algorithmic}[1]
        \State $H \gets \mathrm{MD5}(p)$
        \State $r \gets h_0 \oplus h_1 \oplus \cdots \oplus h_{11}$ \Comment{Equation~\eqref{eq:row}}
        \State $c \gets h_{12} \oplus h_{13} \oplus h_{14} \oplus h_{15}$ \Comment{Equation~\eqref{eq:col}}
        \State $\chi_0 \gets \textsc{Decode}(I[r, 0])$
        \If{$\chi_0 = \mathtt{\$}$}
            \State \textbf{abort:} ``Row $r$ is already occupied.''
        \EndIf
        \State $I'[r, 0] \gets \textsc{Encode}(I[r, 0], \mathtt{\$})$ \Comment{Mark row as occupied}
        \State $t \gets \textsc{CurrentTimestamp}()$
        \State $\pi \gets$ \texttt{*} $\Vert\, w \,\Vert$ \texttt{\#} $\Vert\, t \,\Vert$ \texttt{\#@}
        \For{$i \gets 0$ to $|\pi| - 1$}
            \State $x \gets (c + i) \bmod 256$
            \State $I'[r, x] \gets \textsc{Encode}(I[r, x], \pi[i])$ \Comment{3--3--2 bit split}
        \EndFor
        \State \Return $I'$
    \end{algorithmic}
\end{algorithm}

\subsection{Extraction Algorithm}

The extraction routine reverses the encoding using only the password and the stego-image. It first re-derives $r$ and $c$ from the password's MD5 digest, then reads characters from row $r$ starting at column $c$ until the end marker `\texttt{\#@}' is observed (or a safety limit is reached). The recovered string is then split on `\texttt{\#}' to recover the watermark and the timestamp. Algorithm~\ref{alg:extract} gives the precise procedure.

\begin{algorithm}[H]
    \RaggedRight
    \caption{Row-based LSB watermark extraction}
    \label{alg:extract}
    \algorithmicrequire \quad Stego-image $I'$, password $p$. \\
    \algorithmicensure \quad Recovered watermark $w$ and timestamp $t$.

    \begin{algorithmic}[1]
        \State $H \gets \mathrm{MD5}(p)$
        \State $r \gets h_0 \oplus h_1 \oplus \cdots \oplus h_{11}$
        \State $c \gets h_{12} \oplus h_{13} \oplus h_{14} \oplus h_{15}$
        \State $s \gets$ empty string
        \For{$i \gets 0$ to $255$}
            \State $x \gets (c + i) \bmod 256$
            \State $s \gets s \,\Vert\, \textsc{Decode}(I'[r, x])$
            \If{$s$ ends with \texttt{\#@}}
                \State \textbf{break}
            \EndIf
        \EndFor
        \If{$s$ does not end with \texttt{\#@}}
            \State \textbf{abort:} ``No valid watermark found.''
        \EndIf
        \State $s \gets$ strip trailing \texttt{\#@} and leading \texttt{*} from $s$
        \State $(w, t) \gets$ split $s$ on first occurrence of \texttt{\#}
        \State \Return $(w, t)$
    \end{algorithmic}
\end{algorithm}

\subsection{Worked Example}

Let the password be $p = \texttt{"swordfish"}$, for which the MD5 digest is
\[
    H = \texttt{1cf9d2c2cf301a2c4ce72e9b3f7c1c0a}.
\]
The first twelve bytes of $H$ are $(\mathtt{0x1c}, \mathtt{0xf9}, \mathtt{0xd2}, \mathtt{0xc2}, \mathtt{0xcf}, \mathtt{0x30}, \mathtt{0x1a}, \mathtt{0x2c}, \mathtt{0x4c}, \mathtt{0xe7}, \mathtt{0x2e}, \mathtt{0x9b})$. Their XOR fold equals $\mathtt{0x65} = 101_{10}$, so the watermark is stored in row $r = 101$. The last four bytes $(\mathtt{0x3f}, \mathtt{0x7c}, \mathtt{0x1c}, \mathtt{0x0a})$ XOR-fold to $\mathtt{0x49} = 73_{10}$, so the payload starts at column $c = 73$. With watermark $w = \texttt{"hello"}$ and a representative timestamp $t = \texttt{14/04/32/26/04/2026}$ the constructed payload is
\[
    \pi = \texttt{*hello\#14/04/32/26/04/2026\#@}, \qquad |\pi| = 28.
\]
The 28 characters of $\pi$ are written into pixels $(101, 73), (101, 74), \ldots, (101, 100)$ of the cover image, while pixel $(101, 0)$ receives the `\$' marker. All other $65{,}536 - 29 = 65{,}507$ pixels remain unchanged.

\section{Audio-Steganography Module}
\label{sec:audio_module}

The audio module embeds an XOR-encrypted, timestamped payload into the LSBs of the raw PCM samples of a WAV file. The embedding pipeline is:

\begin{enumerate}[label=(\arabic*)]
    \item Build the payload $\pi = w \,\Vert\, \texttt{\#} \,\Vert\, t \,\Vert\, \texttt{\#@}$.
    \item Encrypt the payload with a one-time-pad-style XOR using the password $p$ as the repeating key:
    \[
        \pi'[i] = \pi[i] \oplus p[i \bmod |p|], \quad i = 0, \ldots, |\pi| - 1.
    \]
    \item Convert each ciphertext character into its 8-bit binary representation and concatenate to form the bitstream $B$.
    \item Append the 16-bit end marker $E = \texttt{1111111111111110}$ to $B$.
    \item Read the WAV samples as a byte array $A$, and replace the LSB of $A[i]$ with the $i$th bit of $B \,\Vert\, E$:
    \[
        A'[i] = (A[i] \,\&\, \mathtt{0xFE}) \mid B_i, \quad i = 0, \ldots, |B \Vert E| - 1.
    \]
    \item Write $A'$ back as a WAV file using the same audio format as the input.
\end{enumerate}

The capacity of the module is therefore $\lfloor (|A| - 16)/8 \rfloor$ ciphertext characters. For example, a one-second 16-bit mono WAV file at $44.1$~kHz contains $88{,}200$ bytes and can store up to $\sim 11{,}023$ characters, which is several orders of magnitude beyond the $16$-character watermark accepted by the GUI.

The extraction routine reverses the pipeline: it scans the LSBs of $A'$ until it finds the end marker $E$, converts the preceding bits back into a ciphertext string, and decrypts the ciphertext with the same XOR key.

This module was kept deliberately simple. WAV was chosen because it exposes raw PCM samples directly, so the embedding rule is easy to implement and easy to explain. This makes the module suitable for demonstration and testing, even though it is not meant to be robust against later compression into formats such as MP3 or AAC.

\section{Text-Steganography Module}
\label{sec:text_module}

The text module is the simplest of the four. Given an input text file $T$, a watermark $w$, a timestamp $t$ and a password $p$, the module:

\begin{enumerate}[label=(\arabic*)]
    \item Builds the payload $\pi = w \,\Vert\, \texttt{\#} \,\Vert\, t \,\Vert\, \texttt{\#@}$.
    \item XOR-encrypts $\pi$ with $p$ to produce the ciphertext $\pi'$.
    \item Appends $\pi'$ to the cover text inside an HTML-style comment:
    \[
        T' = T \,\Vert\, \texttt{\textbackslash n<!--STEGO:} \,\Vert\, \pi' \,\Vert\, \texttt{-->}.
    \]
\end{enumerate}

The cover text is preserved verbatim and the comment is inert with respect to all common text rendering pipelines (HTML, Markdown, plain-text viewers). Extraction locates the substring \texttt{<!--STEGO:}, reads up to the first \texttt{-->} marker and XOR-decrypts the enclosed ciphertext with the password.

In practical terms, the text module behaves more like hidden metadata attachment than classical linguistic steganography. It does not try to alter writing style, change words, or generate natural-language cover text. That choice was intentional. A comment-based design is much easier to verify during testing and much less likely to damage the visible content of the document.

\section{Video-Steganography Module}
\label{sec:video_module}

Video files are harder to handle because any lossy re-encoding step (for example H.264 or H.265) can destroy the LSB values. The video module avoids this by reusing the image module on the first frame and then re-encoding the cover with the lossless FFV1 codec~\cite{niedermayer2013ffv1}. The full pipeline is:

\begin{enumerate}[label=(\arabic*)]
    \item Invoke FFmpeg to extract the first frame of the cover video as a BMP file:
    \begin{quote}
        \texttt{ffmpeg -i input.mp4 -vf "select=eq(n,0)" -vframes 1 -pix\_fmt bgr24 frame0.bmp}
    \end{quote}
    \item Apply Algorithm~\ref{alg:embed} to the top-left $256 \times 256$ region of \texttt{frame0.bmp}, using the same payload format and the same password derivation as the image module.
    \item Use FFmpeg's \texttt{overlay} filter to substitute the watermarked frame back into the cover, re-encoding losslessly with FFV1 into an MKV container so that the LSBs survive.
\end{enumerate}

The output is therefore an \texttt{.mkv} file rather than an \texttt{.mp4} file, which is a deliberate consequence of the requirement for a lossless container. Extraction performs steps (1) and (2) in reverse: it extracts the first frame and runs Algorithm~\ref{alg:extract} on the top-left $256 \times 256$ region.

This approach is a practical compromise. Full video watermarking is a large problem on its own, especially when compression, frame prediction, and temporal consistency are considered. By reusing the image algorithm on a single frame, the project keeps the video module understandable and consistent with the rest of the framework. The trade-off is that robustness remains limited, which is discussed later in the report.

\section{Common Payload Format and Timestamping}
\label{sec:payload_format}

All four modules share the same payload format
\[
    \pi(w, t) = \langle\textit{start}\rangle \,\Vert\, w \,\Vert\, \texttt{\#} \,\Vert\, t \,\Vert\, \texttt{\#@},
\]
where the start marker is `\texttt{*}' for the image and video modules and is empty for the audio and text modules (because in those cases the start of the embedded region is unambiguous). The timestamp $t$ is generated at embedding time and is formatted as
\[
    t = \texttt{HH/MM/ss/dd/MM/yyyy} \enspace ,
\]
which gives a fixed 19-character string. The shared format means that an extracted payload from any modality can be parsed by the same routine: strip the end marker, strip the leading `\texttt{*}' (if any), split on the first `\texttt{\#}', and report the two halves as the watermark and the timestamp respectively.

\section{Computational Complexity}
\label{sec:complexity}

Table~\ref{tab:complexity} summarizes the worst-case time complexity of the embedding and extraction routines for each module. Throughout, $L$ denotes the payload length (at most $39$ for the image module), $N$ the size of the cover, $K$ the password length and $F$ the cost of one FFmpeg invocation.

\begin{table}[H]
    \centering
    \begin{tabular}{|l|c|c|p{4cm}|}
        \hline
        \thead{Module} & \thead{Embed} & \thead{Extract} & \thead{Dominant cost}  \\ \hline
        Image (BMP, $256 \times 256$) & $\Theta(L)$ & $\Theta(L)$ & MD5 of password ($\Theta(K)$) plus $L$ pixel writes/reads.  \\ \hline
        Audio (WAV) & $\Theta(N + L \cdot K)$ & $\Theta(N)$ & Linear scan of audio bytes; XOR is $\Theta(L \cdot K / |p|)$.  \\ \hline
        Text & $\Theta(N + L \cdot K)$ & $\Theta(N)$ & Linear scan to read the file plus XOR encrypt/decrypt.  \\ \hline
        Video (MKV) & $\Theta(F + L)$ & $\Theta(F + L)$ & FFmpeg invocations dominate; pixel work is negligible.  \\ \hline
    \end{tabular}
    \caption{Worst-case time complexity of the four modules in the proposed framework.}
    \label{tab:complexity}
\end{table}

In practice, the image, audio and text modules complete in less than $50$~ms on commodity hardware; the video module is dominated by FFmpeg and takes between one and ten seconds depending on the length and resolution of the cover.

\section{Evaluation Metrics}
\label{sec:metrics}

We use several evaluation measures to fully characterize the proposed system and to compare it against representative single-modality baselines.

\subsection{Perceptual-Quality Metrics}

\textbf{1. Mean Squared Error (MSE).}
\begin{equation}
    \mathrm{MSE} = \frac{1}{N}\sum_{i=1}^{N}(c_i - c'_i)^2,
\end{equation}
where $c_i$ and $c'_i$ are the cover and stego sample values respectively and $N$ is the number of samples (pixels for images, samples for audio). MSE is reported in the same numeric units as the cover; for an 8-bit channel it lies in $[0, 65{,}025]$.

\textbf{2. Peak Signal-to-Noise Ratio (PSNR).}
\begin{equation}
    \mathrm{PSNR} = 10 \cdot \log_{10}\!\left(\frac{255^2}{\mathrm{MSE}}\right) \enspace \text{[dB]}.
\end{equation}
PSNR is the standard imperceptibility measure for spatial-domain watermarking; values above $40$~dB are conventionally considered indistinguishable from the cover.

\textbf{3. Per-Channel Perturbation Bound.} The 3--3--2 bit split guarantees a deterministic upper bound of $\pm 7$ on the red and green channels and $\pm 3$ on the blue channel; this is reported alongside MSE/PSNR as a worst-case quality guarantee.

\subsection{Capacity Metrics}

\textbf{4. Co-Existing-Watermark Capacity.} The number of password-keyed marks that can be embedded into the same cover before the row-occupation marker refuses an embedding. The theoretical maximum for the image module is $256$.

\textbf{5. Per-Cover Payload Length.} The maximum number of ASCII characters that the system embeds per cover (image: $39$; audio: $\sim 11{,}023$ per second of WAV; text: bounded only by the cover size).

\subsection{Correctness Metrics}

\textbf{6. Bit-Exact Recovery.} Whether the recovered $(w, t)$ pair matches the input $(w, t)$ exactly, computed as a binary success/failure indicator across the test corpus.

\textbf{7. Collision-Detection Reliability.} Whether the row-occupation marker correctly refuses an embedding under a colliding password, computed as the fraction of correctly refused attempts out of the total deliberate-collision attempts.

\subsection{Run-Time Metrics}

\textbf{8. Wall-Clock Embed/Extract Time.} The mean elapsed time, in milliseconds, of one embedding or extraction call, averaged over $100$ runs on the same cover. This metric is critical because it determines whether the system is fast enough to be used interactively from a Swing GUI.

\subsection{Comparison Protocol}

We compare the implemented system against four representative single-modality LSB and transform-domain baselines from the literature: Chan and Cheng~\cite{chan2004hiding}, Wu and Tsai~\cite{wu2003steganography}, Cox \emph{et~al.}~\cite{cox1997secure} and Bender \emph{et~al.}~\cite{bender1996techniques}. The comparison is qualitative because the published baselines do not report multi-watermark capacity, password-keyed location selection, or unified payload formats; the qualitative axes are reported in Table~\ref{tab:comparison} of Chapter~\ref{Experimental_results}.

\section{Implementation Details}
\label{sec:implementation}

\subsection{Development Environment}

The system is implemented in Java SE~17 using only the Java standard library for the image, audio and text modules; the video module shells out to the FFmpeg executable for frame extraction, frame substitution and lossless re-encoding. The GUI uses the Java Swing toolkit. The total source-code size is approximately $1{,}500$ lines spread across four classes.

\subsection{Computational Resources}

All experiments described in Chapter~\ref{Experimental_results} run on commodity hardware (multi-core processor, $16$~GB RAM, $512$~GB SSD). The image, audio and text modules complete an embedding or extraction in under $130$~ms; the video module is dominated by FFmpeg and takes one to ten seconds per clip.

\subsection{Reproducibility}

To make the experiments reproducible, the following provisions are implemented:

\begin{itemize}
    \item The same MD5 digest of a given password always yields the same $(r, c)$ pair, so two embeddings of the same watermark with the same password into the same cover produce identical stego-covers (up to the timestamp).
    \item All four modules write a step-by-step log to an in-window console covering the MD5 digest, the XOR folds, the binary encoding of every embedded character, and the per-pixel modification, so that any embedding can be replayed or audited from the log alone.
    \item The cover-file corpus, password set and watermark set used for the experiments are listed explicitly in Section~\ref{sec:datasets}.
    \item The GUI source, the four module sources, and the build/run commands are listed in Section~\ref{sec:hwsw}, so that an undergraduate reader can reproduce the entire pipeline with a single \texttt{javac}/\texttt{java} invocation.
\end{itemize}

\section{Summary}
\label{sec:methodology_summary}

This chapter provided a complete methodology for the proposed multi-modal information-hiding system, including:

\begin{enumerate}
    \item \textbf{System Architecture}: a layered Swing-plus-four-modules pipeline that separates GUI logic from per-modality embedding logic and dispatches user actions through a shared payload format.
    \item \textbf{Image Module}: a row-based LSB scheme with a password-derived $(r, c)$ pair, a 3--3--2 RGB bit split, a row-occupation marker for collision detection, and a self-delimited timestamped payload.
    \item \textbf{Audio Module}: an XOR-encrypted, end-marker-terminated single-bit LSB pipeline on PCM samples.
    \item \textbf{Text Module}: an XOR-encrypted, comment-marker-delimited block appended to the cover text.
    \item \textbf{Video Module}: frame-level reuse of the image module followed by lossless FFV1 re-encoding to preserve the LSBs.
    \item \textbf{Common Payload Format}: a single self-delimited timestamped layout reused across all four carriers.
    \item \textbf{Evaluation}: a four-axis evaluation protocol (perceptual quality, capacity, correctness, run-time) with a qualitative comparison against representative single-modality baselines.
    \item \textbf{Implementation}: a $\sim$1{,}500-line Java SE~17 source base with no external library dependencies for the image/audio/text modules and a single FFmpeg dependency for the video module.
\end{enumerate}

The methodology integrates classical LSB primitives, a cryptographic key-derivation step, and a multi-modal payload contract under a single graphical front-end. The next chapter reports the experimental evaluation that operationalizes the metrics defined in Section~\ref{sec:metrics}.

\textit{\textcolor{red}{Note: The block diagram of Figure~\ref{fig:architecture} satisfies the template's requirement of a clear system-architecture diagram in this chapter.}}
